% LaTeX Package: runcode v2.6
%
% Copyright (C) 2020-2026 by Haim Bar and HaiYing Wang
%
% This file may be distributed and/or modified under the
% conditions of the LaTeX Project Public License, either
% version 1.3c of this license or (at your option) any later
% version.  The latest version of this license is in:
%
%   http://www.latex-project.org/lppl.txt
%
% and version 1.3c or later is part of all distributions of
% LaTeX version 2005/12/01 or later.

\PassOptionsToPackage{colorlinks=true,linkcolor=blue,urlcolor=blue}{hyperref}
\documentclass{ltxdoc}
\usepackage[T1]{fontenc}
\usepackage[utf8]{inputenc}
% hyperref is loaded by ltxdoc via hypdoc; options are set above.
\usepackage{listings}
\usepackage{xcolor}

\lstset{
  basicstyle=\ttfamily\small,
  frame=single,
  breaklines=true,
  columns=flexible,
}

\newcommand{\pkg}[1]{\textsf{#1}}
\newcommand{\opt}[1]{\texttt{#1}}
\newcommand{\file}[1]{\texttt{#1}}

\title{The \pkg{runcode} package}
\author{Haim Bar and HaiYing Wang \\
  \href{mailto:haim.bar@uconn.edu}{haim.bar@uconn.edu},
  \href{mailto:haiying.wang@uconn.edu}{haiying.wang@uconn.edu}}
\date{v2.6, 2026}

\begin{document}
\maketitle

\begin{abstract}
\pkg{runcode} is a \LaTeX{} package that executes programming source
codes (including all command-line tools) from \LaTeX{} and embeds the
results in the resulting PDF file.  Many programming languages can be
used and any command-line executable can be invoked when preparing the
PDF from a \texttt{.tex} file.  \pkg{runcode} is also available on
\href{https://ctan.org/pkg/runcode}{CTAN}.

It is recommended to use this package in server mode together with the
Python \href{https://pypi.org/project/talk2stat/}{talk2stat} package.
Currently, the server mode supports
\href{https://julialang.org/}{Julia},
\href{https://www.mathworks.com/products/matlab.html}{MatLab},
\href{https://www.python.org/}{Python}, and
\href{https://www.r-project.org/}{R}.
More languages will be added.

For more details, usage examples, and troubleshooting, refer to the
package's GitHub repository at
\url{https://github.com/Ossifragus/runcode}.
\end{abstract}

\tableofcontents

%%--------------------------------------------------------------------
\section{Installation}
\label{sec:installation}

You can simply put the \file{runcode.sty} file in your \LaTeX{} project
folder.

The server mode requires the
\href{https://pypi.org/project/talk2stat/}{talk2stat} package.
Some chunk commands (e.g., \cs{showChunk}) also require the
\href{https://pypi.org/project/advance-touch/}{advance-touch} package.
Install them from the command line with:
\begin{lstlisting}[language=bash]
pip3 install talk2stat advance-touch
\end{lstlisting}

\textbf{Note:} \pkg{runcode} requires the \texttt{shell-escape} option
when compiling a \LaTeX{} document:
\begin{lstlisting}[language=bash]
pdflatex --shell-escape yourfile.tex
\end{lstlisting}

When using AUCTeX or a similar IDE, you can avoid setting
\texttt{shell-escape} manually by adding the following magic comment at
the top of your \texttt{.tex} file:
\begin{lstlisting}[language=tex]
% !TEX program = pdflatex --shell-escape
\end{lstlisting}
Replace \texttt{pdflatex} with \texttt{xelatex} or \texttt{lualatex} as
appropriate.

%%--------------------------------------------------------------------
\section{Usage}
\label{sec:usage}

\subsection{Loading the package}

\begin{lstlisting}[language=tex]
\usepackage[options]{runcode}
\end{lstlisting}

Available options:

\begin{description}
  \item[\opt{cache}] Use cached results (do not re-run code).
  \item[\opt{fvextra}] Use the \href{https://ctan.org/pkg/fvextra}{fvextra}
    package to display code.
  \item[\opt{julia}] Start a server for
    \href{https://julialang.org/}{Julia} (requires \pkg{talk2stat}).
  \item[\opt{listings}] Use the
    \href{https://ctan.org/pkg/listings}{listings} package to display code.
  \item[\opt{matlab}] Start a server for
    \href{https://www.mathworks.com/products/matlab.html}{MatLab}
    (requires \pkg{talk2stat}).
  \item[\opt{minted}] Use the
    \href{https://ctan.org/pkg/minted}{minted} package to display code
    (requires \href{https://pygments.org/}{Pygments}).
    This is the default.
  \item[\opt{nominted}] Use \href{https://ctan.org/pkg/fvextra}{fvextra}
    instead of \pkg{minted} (no syntax highlighting; no Pygments required).
  \item[\opt{nohup}] Use \texttt{nohup} when starting a server.
    Some editors (e.g.\ Emacs with AUCTeX) terminate child processes after
    compilation; this option keeps the server alive.
    \textbf{Must be declared before any language option}, e.g.\
    \opt{[nohup,R]} works but \opt{[R,nohup]} does not.
    See Section~\ref{sec:tools} for tools that smooth the \opt{nohup}
    workflow.
  \item[\opt{python}] Start a server for
    \href{https://www.python.org/}{Python} (requires \pkg{talk2stat}).
  \item[\opt{R}] Start a server for \href{https://www.r-project.org/}{R}
    (requires \pkg{talk2stat}).
  \item[\opt{reducedspace}] Reduce the vertical space around output boxes.
  \item[\opt{run}] Force source code to run (override cache).
  \item[\opt{stopserver}] Stop the server(s) when PDF compilation finishes.
\end{description}

\subsection{Output box style}

Output boxes are displayed using
\href{https://ctan.org/pkg/tcolorbox}{tcolorbox} and can be customised
with \cs{tcbset}, e.g.:
\begin{lstlisting}[language=tex]
\tcbset{breakable, colback=red!5!white, colframe=red!75!black}
\end{lstlisting}

\subsection{Code display style}

The style of \emph{source code} blocks (shown by \cs{showCode},
\cs{showChunk}, and the \texttt{\cs{run}LANG\texttt{Chunk}} commands)
depends on the package option used to display code:

\begin{itemize}
  \item \textbf{\pkg{minted}} (default): use \cs{setminted} to control
    language-specific highlighting, e.g.:
\begin{lstlisting}[language=tex]
\setminted[r]{linenos, frame=single, bgcolor=bg, breaklines=true}
\end{lstlisting}
  \item \textbf{\pkg{fvextra}} or \opt{nominted}: use \cs{fvset} to set
    options globally, e.g.:
\begin{lstlisting}[language=tex]
\fvset{frame=single, numbers=left, breaklines=true}
\end{lstlisting}
  \item \textbf{\pkg{listings}}: use \cs{lstset} to configure the style,
    e.g.:
\begin{lstlisting}[language=tex]
\lstset{basicstyle=\ttfamily\small, frame=single, numbers=left}
\end{lstlisting}
\end{itemize}

\subsection{Basic commands}

\DescribeMacro{\runExtCode}
\cs{runExtCode}\marg{prog}\marg{src}\marg{label}\oarg{flag}
runs an external code file.
\begin{itemize}
  \item \meta{prog}: the executable program.
  \item \meta{src}: the source file name.
  \item \meta{label}: the output file name (empty $\Rightarrow$ uses
    counter \texttt{codeOutput}).
  \item \oarg{flag}: optional run control --- omit or empty to use the
    global \texttt{runcode} Boolean; \texttt{run} to force execution;
    \texttt{cache} (or any other value) to use cached results.
\end{itemize}

\textbf{Checksum-based cache invalidation (v2.5):} when neither
\opt{cache} nor \opt{run} is active, \pkg{runcode} computes an MD5
checksum of \meta{src} and compares it with the checksum stored from the
previous run.  The script is re-executed only if the source has changed
or if the output file is missing; otherwise the cached output is reused.
This avoids unnecessary reruns without requiring the user to manage
cache flags manually.

\DescribeMacro{\showCode}
\cs{showCode}\marg{lang}\marg{src}\oarg{first}\oarg{last}
displays source code using \pkg{minted}, \pkg{fvextra}, or \pkg{listings}.
\begin{itemize}
  \item \meta{lang}: programming language.
  \item \meta{src}: source file name.
  \item \oarg{first}/\oarg{last}: first and last line to display
    (optional; defaults to the whole file).
\end{itemize}

\DescribeMacro{\includeOutput}
\cs{includeOutput}\marg{label}\oarg{type}
embeds the output of executed code.
\begin{itemize}
  \item \meta{label}: output file name (must match \meta{label} in
    \cs{runExtCode}; empty $\Rightarrow$ uses counter).
  \item \oarg{type}: \texttt{vbox} (default) = verbatim in a box;
    \texttt{tex} = raw \LaTeX{}; \texttt{inline} = inline text.
\end{itemize}

\DescribeMacro{\inln}
\cs{inln}\marg{prog}\marg{code}\oarg{label}\oarg{type}
runs a short piece of code and displays the result inline.
\begin{itemize}
  \item \meta{prog}: executable or language server.
  \item \meta{code}: source code.
  \item \oarg{label}: output file name (optional).
  \item \oarg{type}: controls display format and caching.  Values:
    \texttt{inline} (default), \texttt{vbox}, \texttt{tex};
    append \texttt{.cache} to any of these
    (e.g.\ \texttt{vbox.cache}) to reuse a cached result,
    re-running only if the output file is absent.
\end{itemize}

\DescribeMacro{\showChunk}
\cs{showChunk}\marg{lang}\marg{src}\marg{id}\oarg{begin}\oarg{end}
displays a labelled chunk of a source file.  Chunks are delimited by
comment lines containing \texttt{label===}\meta{id} (start) and
\texttt{===end} (end) by default.  These delimiters can be overridden
with \oarg{begin} and \oarg{end}.

\DescribeMacro{\writeChunk}
\cs{writeChunk}\marg{lang}\marg{src}\marg{id}\oarg{begin}\oarg{end}
extracts the chunk identified by \meta{id} from \meta{src} and writes it
to \file{generated/}\meta{src}\texttt{-}\meta{id} on disk.  This is
called internally by \cs{showChunk} and \cs{runLANGChunk}, but can also
be used standalone when only the extraction step is needed.

\subsection{Extended commands}

\DescribeMacro{\runCodeIncOut}
\cs{runCodeIncOut}\marg{prog}\marg{src}\oarg{flag}\oarg{label}\oarg{type}
combines \cs{runExtCode} and \cs{includeOutput} in one call.

\subsection{Language-specific shortcuts}

Replace \texttt{LANG} with \texttt{Julia}, \texttt{MatLab},
\texttt{Python}, or \texttt{R} in the following commands.

\DescribeMacro{\runLANG}
\cs{runLANG}\oarg{prog}\marg{src}\marg{label}\oarg{flag}
runs an external \texttt{LANG} code file.  \oarg{prog} is optional and
defaults to the \pkg{talk2stat} \texttt{LANG} server.

\DescribeMacro{\runLANGIncOut}
\cs{runLANGIncOut}\oarg{prog}\marg{src}\oarg{flag}\oarg{label}\oarg{type}
runs a \texttt{LANG} code file and embeds the output.

\DescribeMacro{\inlnLANG}
\cs{inlnLANG}\oarg{prog}\marg{code}\oarg{label}\oarg{type}
runs \texttt{LANG} source code and displays the result inline.
If \meta{code} is wrapped in triple backticks (\texttt{```code```}),
it is sent directly to the server; otherwise it is written to a file
first.

\DescribeMacro{\runLANGChunk}
\cs{runLANGChunk}\oarg{prog}\marg{src}\marg{id}\oarg{flag}\oarg{label}\oarg{type}
extracts the chunk \meta{id} from \meta{src}, runs it, and embeds the
output.
\begin{itemize}
  \item \oarg{label}: output file name.  When omitted, the label is
    derived automatically as \meta{src}\texttt{-}\meta{id}
    (e.g.\ \texttt{code/analysis.R-section2}).
  \item \oarg{type}: output display type (\texttt{vbox}, \texttt{tex},
    \texttt{inline}, or their \texttt{.cache} variants).
\end{itemize}

For example:
\begin{lstlisting}[language=tex]
\runR{code/analysis.R}{result1}
\runRIncOut{code/analysis.R}[][result1]
\runRChunk{code/analysis.R}{section2}
\inlnR{```mean(c(1,2,3,4,5))```}
\end{lstlisting}

\subsubsection{Python batch mode}

\DescribeMacro{\runPythonBatch}
\cs{runPythonBatch}\oarg{src}\oarg{label}
runs a Python file in batch mode (no server).  Requires the
\href{https://pypi.org/project/dill/}{dill} module
(\texttt{pip3 install dill}), which provides session save/restore
so that variables persist across calls.

%%--------------------------------------------------------------------
\section{Auxiliary tools}
\label{sec:tools}

The following scripts are distributed with the package and address
common workflow needs.  Copy them into your project directory alongside
your \texttt{.tex} file.

\subsection{\file{wait\_for\_server.py}}

Polls a \pkg{talk2stat} server until it is ready to accept connections.
This eliminates the server-startup race condition that can occur in
\opt{nohup} mode, where the first \cs{runR} call may arrive before the
server socket is open.

\begin{lstlisting}[language=bash]
python3 wait_for_server.py [LANG [DIR [TIMEOUT]]]
\end{lstlisting}

\begin{itemize}
  \item \texttt{LANG}: language server to wait for --- \texttt{R},
    \texttt{python}, \texttt{julia}, or \texttt{matlab} (default: \texttt{R}).
  \item \texttt{DIR}: working directory of the server (default: \texttt{./}).
  \item \texttt{TIMEOUT}: seconds before giving up (default: 120).
\end{itemize}

\subsection{\file{runcode-Makefile.sample}}

A sample \file{Makefile} for documents compiled in \opt{nohup} server
mode.  Copy it to your project as \file{Makefile} and set the
\texttt{MAIN} and \texttt{LANGS} variables.  The \texttt{make all}
target encodes three best practices:

\begin{enumerate}
  \item Quit any existing server before starting a fresh one, so the
    new server picks up the current \file{.config} settings
    (including any updated \texttt{PIPETIMEOUT}).
  \item Call \file{wait\_for\_rserver.py} to block until the server is
    ready before the first \LaTeX{} pass.
  \item After the draft pass (\texttt{-{}-no-pdf} / \texttt{-draftmode}),
    send a language-appropriate no-op command as a sync barrier
    (\texttt{invisible(NULL)} for R, \texttt{pass} for Python, etc.)
    to ensure all queued scripts have finished writing their output
    files before the final passes begin.
\end{enumerate}

\textbf{PIPETIMEOUT note:} the default timeout in the generated
\file{.config} is 60 seconds.  For long-running scripts, raise it by
editing the \file{.config} file (e.g.\ \texttt{PIPETIMEOUT = 3600})
and restarting the server (\texttt{make stopserver}).

\subsection{\file{consolidate.py}}

Produces a self-contained copy of a \pkg{runcode} project in which all
cached outputs are inlined into the \texttt{.tex} source, so the result
compiles without \pkg{talk2stat} or any language runtime.  This is
useful for submission to Overleaf, a journal, or any recipient who does
not have \pkg{runcode} installed.

\begin{lstlisting}[language=bash]
python3 consolidate.py [--out DIR] [--engine CMD] [--no-compile]
                       [--exclude GLOB] MAIN.tex
\end{lstlisting}

The script copies the project to \file{standalone/} (configurable via
\texttt{-{}-out}), transforms all \texttt{.tex} files by replacing
\cs{includeOutput}, \cs{inlnR}/\cs{inlnPython}/\ldots,
\cs{runRIncOut}, \cs{runRChunk}, etc.\ with the cached content from
\file{generated/}, and substitutes \cs{usepackage}\texttt{\{runcode\}}
with a minimal \pkg{tcolorbox} + \pkg{listings} shim.  Run-only
commands (\cs{runR}, \cs{runPython}, \ldots) are silently removed.
A \texttt{make consolidate} target is provided in
\file{runcode-Makefile.sample}.

%%--------------------------------------------------------------------
\section{Revisions}
\label{sec:revisions}

\begin{itemize}

\item \textbf{v2.6, June 2026:}
  \begin{enumerate}
    \item Checksum-based cache invalidation: source files are
      fingerprinted with MD5; a script is re-executed only when its
      source has changed or when the output file is missing, avoiding
      unnecessary reruns.
    \item New \file{wait\_for\_rserver.py}: eliminates the server-startup
      race condition in \opt{nohup} mode (see Section~\ref{sec:tools}).
    \item New \file{runcode-Makefile.sample}: sample \file{Makefile}
      with server lifecycle management, startup synchronisation, and
      sync barrier (see Section~\ref{sec:tools}).
    \item New \file{consolidate.py}: produces standalone \LaTeX{}
      projects with all cached outputs inlined
      (see Section~\ref{sec:tools}).
  \end{enumerate}

\item \textbf{v2.4.1, June 10, 2026:}
  Bug fixes in \cs{writeChunk}: (1) replaced \texttt{mkdir -p} with
  Python \texttt{os.makedirs} for cross-platform subdirectory creation;
  (2) fixed stream exhaustion when more than 16 \cs{showChunk} calls
  appear in one document (streams are now pre-allocated at package load
  time); (3) removed erroneous \verb|\write18{ad ...}| call in
  \cs{showChunk}.
  Added AUCTeX magic comment to example files.

\item \textbf{v2.4, January 17, 2025:}
  (1) Put file names within \texttt{\$\$} for more robust warnings;
  (2) use the Python \pkg{advance-touch} package to create sub-folders
  for chunk-related commands.

\item \textbf{v2.3, January 12, 2024:}
  Two bug fixes: (1) removed extra space after \cs{inlnX};
  (2) fixed compilation error when an underscore appeared inside R code.

\item \textbf{v2.2, September 8, 2023:}
  Added \cs{showChunk} basic command and \cs{runLANGChunk} commands for
  multiple languages.

\item \textbf{v2.1, June 30, 2023:}
  Detokenize code passed to \cs{inln}.  This prevents special
  \LaTeX{} characters (e.g.\ backslash) from being escaped before
  reaching the interpreter.  Thanks to
  \href{https://github.com/kiryph}{kiryph} for the report.

\item \textbf{v2.0, June 23, 2023:}
  Added \cs{runCodeIncOut}; updated \cs{runExtCode} to override cache
  when output is missing; updated \cs{inln} to accept \texttt{cache}
  in \texttt{Arg4}.  Thanks to
  \href{https://github.com/kiryph}{kiryph} for the suggestions.

\item \textbf{v1.9, June 13, 2023:}
  Updated \cs{inln}: optional \texttt{Arg3} is the output file name,
  optional \texttt{Arg4} is the output type.

\item \textbf{v1.8, January 18, 2023:}
  Added support for the \pkg{listings} package.

\item \textbf{v1.7, August 20, 2022:}
  Renamed the \texttt{tmp/} folder to \texttt{generated/} (CTAN
  requirement); renamed the troubleshooting file.

\item \textbf{v1.6, August 10, 2022:}
  Stop only configured/running servers; new \opt{reducedspace} option;
  default server timeout changed to 60 s; expanded troubleshooting
  document.

\item \textbf{v1.5, July 23, 2022:}
  Removed \texttt{utf8x} option from \pkg{inputenc} (conflict with
  \pkg{hyperref}).

\item \textbf{v1.4, July 18, 2022:} Fixed a bug in cache mode.

\item \textbf{v1.3, May 14, 2022:} Removed hard-coded \pkg{minted}
  options.

\item \textbf{v1.2, May 3, 2022:} Added Python server and batch-mode
  options.

\item \textbf{v1.1, April 17, 2021:} Added \opt{nohup} option;
  improved error handling for missing code files and zero-byte outputs.

\end{itemize}

%%--------------------------------------------------------------------
\section{Contributing}
\label{sec:contributing}

We welcome contributions by opening issues or pull requests on
\href{https://github.com/Ossifragus/runcode}{GitHub}.  Additional
example documents written using \pkg{runcode} are especially
appreciated.

\bigskip\noindent
\textbf{Citing \pkg{runcode}:}
Haim Bar and HaiYing Wang (2021).
\href{https://jds-online.org/journal/JDS/article/103/info}{Reproducible
Science with \LaTeX{}}.
\textit{Journal of Data Science} 19(1), 111--125.
DOI:~10.6339/21-JDS998.

\end{document}
